\chapter{System Requirements Specification and Context Analysis}
\label{chap:srs_context}

\section{Introduction}

This chapter presents the system context analysis of the Horse Racing Tournament Management System (HRTMS). The objective of this chapter is to identify the system stakeholders, define the system boundary, and describe how external actors interact with the system.

The context diagram provides a high-level view of the system and illustrates the information exchanged between the system and its external entities.

\section{Stakeholder Analysis}

The Horse Racing Tournament Management System involves several stakeholders who interact with the system for different purposes.

\begin{table}[H]
	\centering
	\caption{System Stakeholders}
	\begin{tabular}{|p{4cm}|p{9cm}|}
		\hline
		\textbf{Stakeholder} & \textbf{Description} \\
		\hline
		Admin & Responsible for managing tournaments, users, race schedules, registrations, referees, and publishing results. \\
		\hline
		Horse Owner & Owns and manages horses, registers horses for tournaments, invites jockeys, and follows race results. \\
		\hline
		Jockey & Participates in races, accepts or rejects invitations, and views assigned races. \\
		\hline
		Race Referee & Inspects races, records violations, inputs race results, and confirms competition outcomes. \\
		\hline
		Spectator & Views tournament information, race schedules, results, rankings, and prediction information. \\
		\hline
	\end{tabular}
\end{table}

\section{System Context Diagram}

Figure~\ref{fig:context_diagram} illustrates the overall context of the Horse Racing Tournament Management System and its interactions with external actors.

\begin{figure}[H]
	\centering
	\includegraphics[width=\textwidth]{images/chuong3/context_diagram.png}
	\caption{Context Diagram of Horse Racing Tournament Management System}
	\label{fig:context_diagram}
\end{figure}

The context diagram represents the system as a single process and shows the information exchanged between the system and the surrounding actors.

\section{Actor Interaction Description}

\subsection{Admin}

The Admin is responsible for overall system administration and tournament management.

The Admin interacts with the system through the following functions:

\begin{itemize}
	\item Register and login.
	\item Manage user permissions.
	\item Create and manage tournaments.
	\item Approve or reject race registrations.
	\item Create race schedules.
	\item Assign referees to races.
	\item Publish race results and rankings.
	\item Manage user, registration, race, and result information.
\end{itemize}

\subsection{Horse Owner}

Horse Owners manage horses participating in tournaments.

The Horse Owner can:

\begin{itemize}
	\item Register and login.
	\item Manage horse information.
	\item Register horses for tournaments.
	\item Invite jockeys.
	\item View registration status.
	\item View race schedules.
	\item View race results.
	\item Track horse rankings.
\end{itemize}

\subsection{Jockey}

Jockeys participate in races by riding horses assigned by Horse Owners.

The Jockey can:

\begin{itemize}
	\item Register and login.
	\item Manage profile information.
	\item Accept or reject invitations.
	\item View assigned race schedules.
	\item View horse and race information.
	\item Monitor participation status.
\end{itemize}

\subsection{Race Referee}

Race Referees supervise race activities and ensure fair competition.

The Race Referee can:

\begin{itemize}
	\item Inspect races before and after competition.
	\item Record race violations.
	\item Input race results.
	\item Confirm race outcomes.
	\item Access assigned race information.
	\item View race schedules.
\end{itemize}

\subsection{Spectator}

Spectators are public users who monitor tournament activities.

The Spectator can:

\begin{itemize}
	\item View tournament information.
	\item View race schedules.
	\item View race results.
	\item View ranking tables.
	\item Submit race predictions.
	\item View prediction results.
\end{itemize}

\section{System Boundary Analysis}

The system boundary defines which functionalities belong to the Horse Racing Tournament Management System and which entities are external to the system.

The following components are located inside the system boundary:

\begin{itemize}
	\item Authentication and Authorization
	\item User Management
	\item Horse Management
	\item Jockey Management
	\item Tournament Management
	\item Registration Management
	\item Race Scheduling
	\item Referee Management
	\item Result Management
	\item Ranking Management
	\item Spectator Services
\end{itemize}

The following entities are located outside the system boundary:

\begin{itemize}
	\item Admin
	\item Horse Owner
	\item Jockey
	\item Race Referee
	\item Spectator
\end{itemize}

The system acts as a centralized platform that facilitates communication and information exchange among all stakeholders.

\section{Data Flow Summary}

Based on the context diagram, the primary information flows within the system can be summarized as follows:

\begin{itemize}
	\item Horse Owners submit horse registrations and receive approval status.
	\item Horse Owners send jockey invitations and receive responses.
	\item Admin manages tournaments, schedules, and approvals.
	\item Referees submit race inspections, violations, and race results.
	\item The system updates rankings after result confirmation.
	\item Spectators access public information including schedules, results, and rankings.
\end{itemize}

These information exchanges form the foundation of the business processes implemented within the Horse Racing Tournament Management System.

\section{Summary}

This chapter presented the context analysis of the Horse Racing Tournament Management System. The stakeholders, system boundary, actor interactions, and major information flows were identified and analyzed.

The context diagram provides a high-level understanding of how the system interacts with external actors and serves as the foundation for the detailed requirement analysis and UML modeling presented in subsequent chapters.
